% October 13, 2026 build (agent draft): introduction written from the Oct 12
% 7.9 deck (frames Motivation, This Paper, Findings, Related Literature) in
% the register of Boar, Gorea and Midrigan (2022). Literature descriptions
% follow notes/literature_slide_20261011/paper_notes_for_tommaso.md; where
% that note says only the abstract was read, the description here is limited
% to the abstract's claim. Numbers: 7p9167_v6/table.md (7.9, production).
\section{Introduction}
\label{sec:intro}

Fertility in the United States has fallen from about 2.1 births per woman in
2007 to 1.62 in 2023. Housing costs are often blamed, and the remedies proposed
range from easier mortgage credit to building more homes. Both remedies rest
on the observation that children need space. Families with children live in
larger homes, and most large homes are sold rather than rented, so a growing
family that wants more space must usually buy it, with a down payment, at an
age when its savings are smallest and its earnings most uncertain. Whether
cheaper credit or more housing raises fertility, and by how much, depends on
which part of this cost holds families back: the one-time cost of getting
into a larger home, or the recurring cost of occupying it.

This paper answers the question with a life-cycle model in which households
choose when to have children, how much space to occupy, and whether to rent or
own it. Three features of the environment carry the analysis. First, a child
raises the space a household needs, so a birth commits the household to
paying for more space in every period the child is at home. Second, rental
units are capped in size: the largest homes can only be bought, and buying
requires a down payment. Third, earnings are subject to uninsurable
idiosyncratic risk, so the commitment a child entails is costly to bear for a
household without savings. The house price clears a single market for owned
and rented rooms, competitive landlords set rents equal to the user cost of
housing, and housing supply responds to the price. The economy is an
overlapping-generations economy whose population is determined by fertility,
so the price of space today depends on how many children were born a
generation ago.

I calibrate the model to the United States in 2007, when fertility was at
replacement, to match fertility, housing and wealth moments: childlessness,
family size and the timing of first births, mean rooms, the ownership rate,
the gap in ownership between recent parents and households without children,
the rooms households add around the first birth, and the ratios of wealth and
bequests to earnings. I then hold the household parameters fixed and
attribute the decline in fertility since 2007 to changes in the preference for
children, chosen to match period fertility through 2023. The 2023 economy is
thus a point on a transition to a smaller population, along which smaller
cohorts enter adulthood and the demand for space falls. Every policy I
consider starts from that economy and is compared with the path the economy
would follow without it.

The main results are as follows. Earnings risk matters for fertility. A
mean-preserving reduction in the standard deviation of earnings shocks from
0.48 to 0.35 per four years raises the completed fertility of the cohort
entering adulthood in 2023 by 16 percent. Housing costs also matter. An
outward shift of the housing supply curve by 10 percent lowers the house price
by 8.9 percent on impact and raises the completed fertility of the same cohort
by 4.8 percent, and of the cohort entering adulthood twenty years later by 3.8
percent. Easier mortgage credit mainly changes the timing of births. Allowing
households to borrow the full value of the home raises births by 1.9 percent
in 2023--2026, but leaves the completed fertility of the 2023 cohort 0.6
percent lower: credit lets young households buy earlier, and the homes they
buy are small. The same credit restricted to households with a child at home
raises completed fertility by 6.3 percent, and by 11.5 percent when earnings
risk is lower.

These results have a common explanation. Because rent equals the user cost of
owning and households and landlords borrow at the same rate, a room costs the
same each period whether it is rented or owned. Mortgage credit changes the
equity a household needs to buy, not the per-period cost of the space a child
requires. A larger housing supply lowers that cost for every family and every
period; lower earnings risk makes it easier to bear. A comparison makes the
point sharply. A supply expansion of 6.3 percent, sized so that it raises
births on impact as much as full mortgage credit does, raises the completed
fertility of the 2023 cohort by 3.1 percent, against $-0.6$ percent for
credit. Credit tied to children is the exception that proves the rule: it
rewards the birth itself with access to ownership.

The results come with three qualifications that the paper discusses in
detail. First, the reduction in earnings risk also compresses the distribution
of earnings over time, and births rise steeply with earnings at the bottom, so
part of the 16 percent may reflect less inequality rather than less
uncertainty. Second, credit restricted to parents is modelled without default
or mortgage insurance, so its cost is not priced. Third, the 2007--2023
transition is driven only by preference shocks: the model overstates
childlessness in 2023, misses the rise in the age at first birth, and
predicts a falling house price where the data show a rising one. The
quantitative results in this draft are at the current calibration and will be
refreshed when the 2007--2023 preference path is re-estimated as a gradual
decline.

\paragraph{Related literature.} This paper contributes to three strands of
work. The first studies how housing costs affect fertility.
\citet{simonTamuraHigherRents2009} show that rent per room is negatively
related to fertility across US cities, and \citet{stromHousingFirstBirths2010}
that dwelling size is the housing feature most strongly tied to first births
in Sweden. Dettling and Kearney (2014) find that higher house prices raise
births among owners and lower them among renters, and Lovenheim and Mumford
(2013) that housing wealth raises owners' births. On credit,
\citet{hacamoBabiesMortgageDeregulation2021} finds that young households
exposed to mortgage deregulation became more likely to buy a home and have a
child, \citet{dettlingKearneyModernMortgage2025} attribute part of the baby
boom to federal mortgage programs that lowered down payments, and
\citet{vandoornikFazioRamadoraiSkrastinsHousingFertility2025} find that
winning a Brazilian housing-credit lottery raises births. Cumming and Dettling
(2024) show that lower mortgage payments raise births among existing owners.
Structural work includes Khorunzhina and Miller (2022), who estimate a model of
births and the first home purchase with exogenous prices, and
\citet{couillardBuildBabyBuild2025}, who studies rents, living arrangements and
fertility in a sorting model. Relative to this work, I study fertility in an
equilibrium model with saving, earnings risk, mortgage finance, tenure and a
house price, which separates the recurring cost of space from the cost of
buying. The model reconciles the evidence on credit with a small effect of
untargeted credit on completed fertility: credit raises births where it is
tied to children or reaches a market whose supply expands, as in the episode
studied by \citet{dettlingKearneyModernMortgage2025}.

The second strand studies how earnings risk affects fertility. Fraser (2001)
shows that prudent parents facing an irreversible commitment want fewer
children under more risk. Sommer (2016) quantifies the effect in a life-cycle
model in which children require a minimum of spending, and de la Croix and
Pommeret (2021) model the first birth as an irreversible investment set
against a biological clock. Ejrn\ae s and J\o rgensen (2020) and Nakakuni (2026)
find that removing earnings risk raises fertility by 6 to 16 percent;
empirical work links job loss and job insecurity to lower fertility (Del Bono,
Weber and Winter-Ebmer, 2012; Clark and Lepinteur, 2022). I make housing the
commitment, so that earnings risk and the cost of space interact, and compare
risk with housing policies in the same equilibrium.

The third strand builds quantitative models of housing over the life cycle,
with tenure choice, down payments, earnings risk and equilibrium prices
(Fisher and Gervais, 2011; Halket and Vasudev, 2014; Kaplan, Mitman and
Violante, 2020; Boar, Gorea and Midrigan, 2022; Paz-Pardo, 2024; Greaney et al.,
2025; Greenwald and Guren, 2025). In these models household size is
exogenous. I make it a choice and tie the demand for space to children. Fisher
and Gervais (2011) find that earnings risk, not the down payment, explains the
fall in young ownership; the analogue here is that earnings risk, not the down
payment, governs completed fertility.

The rest of the paper is organized as follows. Section~\ref{sec:empirical-evidence}
documents facts about housing and the first birth that shape the model.
Section~\ref{sec:model} presents the model, and Section~\ref{sec:quant} its
quantification and the transition from 2007 to 2023.
Section~\ref{sec:results} describes how a first child changes household
choices and why earnings risk matters. Section~\ref{sec:policy} studies
housing supply, mortgage credit and credit for parents from the 2023 economy.
Section~\ref{sec:conclusion} concludes.
